../cictikz/src/cictikz/data/tex/cictikz_preamble.tex

% The same switched capacitor as l05_fund1, rotated: C1 now floats between the
% phi_1 switch and V_O, and phi_2 shorts it out instead of pulling it to
% ground. The charge difference per cycle is the same, so Z_I is again
% 1/(C1 f_phi) -- which is the point the slide makes.
%
% Geometry is shared with l05_fund3 -- same V_I port at -\grid*1.5, same phi_1
% switch from 0 to \grid, same C1 span from \grid to \grid*2, same V_O column
% at \grid*3, same ground rail. Keep them in step if either is re-laid-out.
%
% C1's label goes underneath here rather than above, because the phi_2 switch
% sits over it.

\begin{circuitikz}[american, thick, transform shape, circuit ee IEC]

  ../cictikz/src/cictikz/data/tex/cictikz_lib.tex

  \def\xin{-\grid*1.5}   % V_I port
  \def\xa{0}             % phi_1 switch, left end
  \def\xna{\grid}        % left plate of C1
  \def\xnb{\grid*2}      % right plate of C1
  \def\xout{\grid*3}     % V_O column
  \def\ysig{0}           % signal rail
  \def\yphi{\grid}       % the shorting switch, one grid above the signal
  \def\ygnd{-2.8}        % ground rail

  % ---- Input ----
  \draw (\xin,\ysig) to [short,o-] (\xa,\ysig);
  \node[anchor=east] at (\xin-0.15,\ysig) {$V_{I}$};

  \draw[->] (\xin+0.1,\ysig-0.55) -- ++(1.0,0);
  \node[anchor=north] at (\xin+0.6,\ysig-0.65) {$Z_{I} = \;?$};

  % ---- phi_1 switch ----
  \draw (\xa,\ysig) to [Tnmos, n=MS1] (\xna,\ysig);
  \draw (MS1.gate) to [short,-o] ++(0,0.9)
        node [anchor=south, inner sep=5pt] {$\phi_{1}$};

  % ---- C1 floating between the two nodes ----
  \draw (\xna,\ysig) \hcapacitor{};
  \draw (capEnd) -- (\xnb,\ysig);
  \node[anchor=north] at (\xna+\grid/2,\ysig-0.3) {$C_{1}$};

  % ---- phi_2 shorts C1 out ----
  \draw (\xna,\ysig) -- (\xna,\yphi);
  \draw (\xna,\yphi) to [Tnmos, n=MS2] (\xnb,\yphi);
  \draw (\xnb,\yphi) -- (\xnb,\ysig);
  \draw (MS2.gate) to [short,-o] ++(0,0.9)
        node [anchor=south, inner sep=5pt] {$\phi_{2}$};

  \fill (\xna,\ysig) circle (0.06);
  \fill (\xnb,\ysig) circle (0.06);

  % ---- Output, held by a source as in the original ----
  \draw (\xnb,\ysig) -- (\xout,\ysig);
  \draw (\xout,\ysig) to [short,*-o] ++(0.6,0)
        node [anchor=west, inner sep=5pt] {$V_{O}$};

  \draw (\xout,\ysig) -- ++(0,-0.4);
  \draw (\xout,\ysig-0.4) to [V] (\xout,\ygnd+0.4);
  \draw (\xout,\ygnd+0.4) -- (\xout,\ygnd);
  \draw (\xout,\ygnd) \vground;

\end{circuitikz}
\end{document}
